\section{Verification purpose}
This case verifies dead-fuel moisture damping, heat of preignition, and extinction for the dead-dominated static fuels FBFM3 and TL6.  Moisture is varied as a fraction of each model's configured dead extinction moisture.  It can expose percent/fraction errors, incorrect polynomial evaluation, wrong extinction comparison, or omission of moisture from the heat sink.  The verification question is whether reaction intensity and ROS follow the independent calculation through extinction for every required variant.

\section{Mathematical formulation and numerical implementation}
For $r_M=M_f/M_x$, Rothermel moisture damping is
\begin{equation}\eta_M=\max\{0,\min[1,1-2.59r_M+5.11r_M^2-3.52r_M^3]\}.\end{equation}
Moisture also increases $Q_{ig}=250+1116M_f$.  Thus it reduces the reaction source and increases the heat sink.  The equations follow Andrews (2018), RMRS-GTR-371.  ELMFIRE receives percent rasters; the oracle uses fractions.

\section{Derivation of expected behavior}
For each fuel, the 1-h, 10-h, and 100-h moistures are equal, so their weighted dead moisture is the selected ratio times $M_x$.  The damping polynomial and therefore dead reaction intensity reach zero at and above $r_M=1$.  Because these variants have no live load contribution, $I_R=R_0=0$ there.  Below extinction, the independent fuel weighting and heat sink give the exact expected curve.

\section{Verification metrics and acceptance criteria}
Maximum relative errors are 0.02 for direct ROS and 0.01 for reaction intensity.  Variants predicted to spread at least \SI{0.10}{m.min^{-1}} additionally require TOA-derived ROS error at most 0.05 and regression $R^2\ge0.995$; slower variants cannot cross three \SI{10}{m} cells within the capped practical duration and are decided by direct ROS and intensity.  Extinguished variants require zero-valued direct fields; TOA is not used to invent a velocity where no propagation occurs.  All 16 variants are mandatory.

\section{Simulation configuration and predicted outcome}

The finite-strip PHI raster is a signed-distance field clipped to ELMFIRE's supported $[-1,1]$ interval, preserving both an ignited interior and a resolved front gradient. For every positive oracle ROS, \texttt{SIMULATION\_DT} and \texttt{SIMULATION\_DTMAX} are selected so the expected front advances 5\% of one cell per step (CFL 0.05, below the configured 0.2 limit). This keeps the PHI change above ELMFIRE's fixed stall scale even if the implementation produces only a small fraction of the reference speed. A zero-ROS variant retains a 5 s step and is evaluated from direct fields without requiring TOA propagation. The analytically selected stop time is rounded upward to a whole second because the current ELMFIRE time loop stores the duration as an integer; this preserves the intended duration and guarantees the terminal output dump. The evaluator always requires uniquely named rasters tied to the one final dump record. A TOA-required variant must reach the \texttt{SIMULATION\_TSTOP} in its generated namelist; a direct-only variant may terminate early because lack of propagation is part of the state being tested.
The sweep ratios are 0.10, 0.25, 0.50, 0.70, 0.85, 0.95, 1.00, and 1.10.  Domains use \SI{10}{m} cells and a finite ignition strip with a planar east face in the interior measurement corridor; the strip terminates inside the lateral boundaries.  Wind and slope are zero, live moisture is fixed, and no coupled fire processes are enabled.  Intensity and ROS should decrease toward zero and remain zero beyond extinction.

\section{Actual simulation results}
The generated record contains \CompletedVariantCount{} of \RequiredVariantCount{} required variants and reports \textbf{\VerificationStatus}.
\ifnum\CompletedVariantCount>0
\IfFileExists{../figures/sweep_response.pdf}{\begin{figure}[H]\centering\includegraphics[width=0.95\textwidth]{sweep_response.pdf}\caption{Hand calculations (open diamonds, solid lines) and ELMFIRE predictions (filled circles, dashed lines) for the dead-moisture response through extinction. Each fuel-model pair uses the same color.}\end{figure}}{}
\fi

\section{Calculated metrics and verification decision}
\begin{center}\begin{tabular}{@{}llll@{}}\toprule Metric & Limit & Calculated & Status\\\midrule\MetricRows\bottomrule\end{tabular}\end{center}
The authoritative decision is \textbf{\VerificationStatus}.

\section{Technical interpretation}
A horizontal shift of the extinction point indicates an $M_x$ or units error.  Correct intensity with wrong ROS suggests heat-sink or flux-ratio handling.  Nonzero propagation beyond extinction indicates an unintended live contribution, adjustment, or stale output.

\section{Scope and limitations}
The case varies all three dead moisture classes together.  It does not identify an isolated class-indexing defect or validate moisture observations.  Static live moisture is addressed separately by CASE38.
